\begin{lemma}
Let \(U_0,V_0\subseteq\operatorname{Fin}n\) be fixed disjoint sets of
original vertices.  Fix a blue base graph \(G_B\), a blue row map
\(\rho:\operatorname{Fin}n\to\operatorname{Fin}m\), and an exposure
\(\eta:U_0\to\operatorname{Fin}m\).  Assume the blue rows used by \(U_0\) and
\(V_0\) are disjoint:
\[
u\in U_0,\ v\in V_0 \quad\Longrightarrow\quad \rho(u)\ne\rho(v).
\]
Put \(S(\eta)=\eta(U_0)\).  Conditional on \(G_B\), on the full blue row map
\(\rho\), and on the red-coordinate exposure \(\eta\) on \(U_0\), the
probability that at least \(t\) vertices of \(V_0\) have red coordinate in
\(S(\eta)\) is at most
\[
\binom{|V_0|}{t}\left(\frac{|S(\eta)|}{m}\right)^t .
\]
\end{lemma}
\begin{proof}
Let
\[
C_0=\{\omega:\omega.2.1=G_B\text{ and }(\pi_B^\omega(v))=\rho(v)
\text{ for every }v\},
\]
where \(\pi_B^\omega(v)\) denotes the second coordinate of the injection in
the sample \(\omega\).  Let
\[
H_\eta=\{\omega:(\pi_R^\omega(u))=\eta(u)\text{ for every }u\in U_0\},
\qquad C_\eta=C_0\cap H_\eta .
\]
If \(\operatorname{TwoBiteEventProbability}(C_\eta)=0\), then the conditional
probability in \(\noderef{TwoBiteConditionalProbability}\) is defined to be
\(0\), and the desired inequality is immediate.

Assume now that \(C_\eta\) has nonzero mass.  Expand the conditional
probability over the finite two-bite sample space.  In the conditioning cell
\(C_0\), the blue graph and the blue row map are fixed.  By
\(\noderef{TwoBiteSampleWeight}\), all remaining sample weights factor into a
common factor depending on \(G_B\), a red-graph weight, and the uniform
injection weight.  The sum over red graphs is independent of the injection
coordinates, and \(\noderef{GnpGraphWeightSumOne}\) makes this red-graph
sum equal to \(1\).  Thus, after the common blue-graph and
uniform-injection factors are canceled, only the finite count of compatible
injections remains.  Those injection counts split over the blue rows.

Now pass from sample weights to injection counts.  Let \(D_\eta\) be the set
of injections \(\phi:\operatorname{Fin}n\hookrightarrow
\operatorname{Fin}m\times\operatorname{Fin}m\) satisfying
\[
(\phi(v))_2=\rho(v)\quad\text{for every }v,
\qquad
(\phi(u))_1=\eta(u)\quad\text{for every }u\in U_0.
\]
For a fixed set \(T\subseteq V_0\), let \(D_\eta(T)\subseteq D_\eta\) be the
subset satisfying \((\phi(v))_1\in S(\eta)\) for every \(v\in T\).  Since the
conditioning mass is nonzero, the weighted denominator for \(C_\eta\) is
nonzero; after the common blue-graph, red-graph-sum, and uniform-injection
factors are canceled, this implies \(|D_\eta|>0\).  Therefore conditional
probabilities for events depending only on the injection are ratios
\[
\Pr(D_\eta(T)\mid C_\eta)=\frac{|D_\eta(T)|}{|D_\eta|}.
\]

Decompose these two finite sets by blue rows.  For each row
\(i\in\operatorname{Fin}m\), put
\[
R_i=\{v:\rho(v)=i\},\qquad A_i=U_0\cap R_i,\qquad
B_i=T\cap R_i.
\]
The denominator row factor \(d_i\) is the number of injective red-coordinate
assignments on \(R_i\) that agree with \(\eta\) on \(A_i\).  The numerator
row factor \(d_i(T)\) imposes in addition that the vertices of \(B_i\) land
in \(S(\eta)\).  The global counts factor as
\[
|D_\eta|=\prod_i d_i,\qquad |D_\eta(T)|=\prod_i d_i(T),
\]
because choices in different blue rows are independent and pairs with
different second coordinate are automatically distinct.

If \(B_i=\varnothing\), then the extra \(T\)-condition imposes nothing in row
\(i\), so \(d_i(T)=d_i\).  If \(B_i\ne\varnothing\), then row-disjointness
implies \(A_i=\varnothing\): a row containing a vertex of \(T\subseteq V_0\)
cannot contain a vertex of \(U_0\).  Thus the exposure \(\eta\) imposes no
condition in this row.  We now count this row directly.

Write
\[
q_i=|R_i|,\qquad h_i=|B_i|,\qquad s=|S(\eta)|.
\]
The denominator row count is the number of injective red-coordinate
assignments on the \(q_i\) vertices of \(R_i\), namely the falling factorial
\((m)_{q_i}=m(m-1)\cdots(m-q_i+1)\).  Since the present conditioning cell has
nonzero denominator, each row factor \(d_i\) is nonzero; hence in the present
row \(q_i\le m\).  If \(h_i>s\), then the numerator row count is \(0\), since
the \(h_i\) distinct vertices in \(B_i\) would have to receive distinct
red coordinates inside the \(s\)-element set \(S(\eta)\).  In that case the
desired row bound is immediate.

It remains to consider \(h_i\le s\).  First choose injectively the
red coordinates of the \(h_i\) constrained vertices from \(S(\eta)\), and
then choose injectively the remaining \(q_i-h_i\) red coordinates from the
unused coordinates.  This gives
\[
d_i(T)=(s)_{h_i}(m-h_i)_{q_i-h_i},
\qquad
d_i=(m)_{q_i}=(m)_{h_i}(m-h_i)_{q_i-h_i}.
\]
The common factor \((m-h_i)_{q_i-h_i}\) cancels, so
\[
\frac{d_i(T)}{d_i}
=\frac{(s)_{h_i}}{(m)_{h_i}}
=\prod_{\ell=0}^{h_i-1}\frac{s-\ell}{m-\ell}.
\]
For \(0\le\ell<h_i\le s\le m\), we have
\[
\frac{s-\ell}{m-\ell}\le \frac{s}{m},
\]
because after multiplying by the positive denominators this is equivalent
to \(m(s-\ell)\le s(m-\ell)\), i.e. \(\ell s\le \ell m\).  Multiplying these
\(h_i\) factor bounds gives the row estimate
\[
\frac{d_i(T)}{d_i}\le
\left(\frac{|S(\eta)|}{m}\right)^{|B_i|}.
\]
All denominator factors \(d_i\) are nonzero in the present nonzero cell; if
some \(d_i=0\), then \(|D_\eta|=0\), contrary to \(|D_\eta|>0\).
Multiplying the row ratios and using that the sets \(B_i=T\cap R_i\)
partition \(T\), we get
\[
\frac{|D_\eta(T)|}{|D_\eta|}
\le
\prod_i\left(\frac{|S(\eta)|}{m}\right)^{|B_i|}
=
\left(\frac{|S(\eta)|}{m}\right)^{|T|}.
\]
Equivalently,
\[
\Pr(\forall v\in T,\ \pi_R^\omega(v)\in S(\eta)\mid C_\eta)
\le \left(\frac{|S(\eta)|}{m}\right)^{|T|}.
\]

Let
\[
W_\eta=|\{v\in V_0:\pi_R^\omega(v)\in S(\eta)\}|.
\]
If \(W_\eta\ge t\), then there is a \(t\)-element subset
\(T\subseteq V_0\) such that every \(v\in T\) has red coordinate in
\(S(\eta)\).  Taking the finite union over all such \(T\) gives
\[
\Pr(W_\eta\ge t\mid G_B,\rho,\eta)
\le \binom{|V_0|}{t}\left(\frac{|S(\eta)|}{m}\right)^t.
\]
The same estimate is trivial when \(t>|V_0|\), since there are no
\(t\)-element subsets of \(V_0\).  This proves the claimed bound.
\end{proof}
